\documentclass[runningheads]{llncs}
% 
\usepackage[T1]{fontenc}
\usepackage{graphicx}
\usepackage{amsmath}
\usepackage{amssymb}
\usepackage{siunitx}
\usepackage{algorithm}
\usepackage{algpseudocode}
\usepackage{booktabs}
\usepackage[hidelinks]{hyperref}
\usepackage{tabularx}

\usepackage{makecell}
\usepackage{subcaption}

\begin{document}

\title{Transformer-Based Spectral Unmixing for Lunar Mineral Mapping}
\author{
Nishant\inst{1} \and
Bindu Verma\inst{1} \and
Kapil Sharma\inst{1}
}

\authorrunning{Nishant et al.}

\institute{
Department of Information Technology, Delhi Technological University, Delhi, India\\
\email{nishantsinghchauhan0010@gmail.com, bindu.cvision@gmail.com, kapil@ieee.org}
}
\maketitle              
%

\begin{abstract}
Lunar mineral mapping remains challenging because of spectral dimensionality, mixing, nonlinear variability, and long-range dependencies. This study presents a one-dimensional Vision Transformer (ViT) for spectral unmixing and mineral abundance estimation using Chandrayaan-2 Imaging Infrared Spectrometer (IIRS) observations. Each spectrum is represented as spectral patches, and self-attention learns long-range relationships. The dual-head architecture reconstructs spectra and predicts non-negative, sum-to-one mineral fractions through a Softmax abundance head; chunk-wise inference enables full-scene mapping. On held-out IIRS observations, the model achieved an RMSE of 0.03637, MAE of 0.01420, and SAD of 0.02314 for spectral reconstruction, alongside an abundance-sum mean of 1.00000 and zero negative-abundance ratio. Full-scene inference produced predicted mineral distributions, with mare basalt, pyroxene, and spinel accounting for 63.61\%, 16.90\%, and 14.92\% of mapped pixels, respectively. These results demonstrate the potential of one-dimensional transformer-based spectral learning for lunar mineral mapping.




\keywords{Hyperspectral imaging \and Vision Transformer \and Spectral
unmixing \and Lunar mineral mapping \and Chandrayaan-2 IIRS \and
Mineral abundance estimation}


\end{abstract}


\section{Introduction}

Hyperspectral imaging (HSI) is an important remote sensing technology
for planetary surface analysis, providing continuous spectral information
across hundreds of wavelength bands. In lunar exploration, HSI supports
the characterization of surface mineralogy and geological composition.
The Imaging Infrared Spectrometer (IIRS) onboard Chandrayaan-2 acquires
256 spectral bands spanning approximately $0.8$--$5.0~\si{\micro\metre}$,
providing diagnostic spectral information for lunar mineral
identification \cite{chowdhury2020iirs,chauhan2021hydration}.

Despite this capability, lunar hyperspectral analysis remains
challenging owing to high spectral dimensionality, spectral
redundancy, instrumental noise, illumination variability, and spectral
mixing. Deep learning has improved hyperspectral
representation learning through nonlinear feature extraction
\cite{zhang2016deeplearning,paoletti2019review}. However, CNN-based
methods primarily capture local spectral patterns and may have limited
capability to model long-range dependencies across spectral bands.

Transformer-based methods address this limitation through self-attention
and global spectral representation learning. SpectralFormer,
spatial-spectral transformers, and dual attention transformers have
demonstrated the effectiveness of transformer-based learning for
hyperspectral analysis \cite{hong2022spectralformer,he2021sst,shu2024datn}.
Vision Transformers (ViTs) represent inputs as token sequences
processed through multi-head self-attention
\cite{dosovitskiy2021vit}. Nevertheless, most existing methods
primarily address supervised hyperspectral classification, with
comparatively limited exploration of transformer-based spectral
unmixing for quantitative mineral abundance estimation.

An important distinction exists between spectral reconstruction and
mineral abundance estimation. Reconstruction measures how closely
predicted spectra approximate observed reflectance, whereas abundance
estimation determines the fractional contribution of individual mineral
components. Consequently, low reconstruction error alone does not
establish geologically reliable mineral abundance estimates.

To address these limitations, this study proposes a one-dimensional
Vision Transformer framework for spectral unmixing and lunar mineral
mapping using Chandrayaan-2 IIRS observations. Each 256-band spectrum is
divided into contiguous 16-band spectral patches and projected onto a
token sequence. A pre-layer normalization ViT with six encoder layers,
256-dimensional embeddings, and eight attention heads models local
spectral characteristics and long-range dependencies
\cite{dosovitskiy2021vit,xiong2020layernorm}. The shared latent
representation supports two complementary heads for spectral
reconstruction and mineral abundance prediction, with the latter
subject to non-negativity and sum-to-one constraints.

The major contributions of this work are summarized as follows:

\begin{enumerate}

\item \textbf{A one-dimensional Vision Transformer framework} was
proposed to learn long-range spectral dependencies from Chandrayaan-2
IIRS observations using spectral patch tokenization.

\item \textbf{A dual-head learning strategy} jointly performs spectral
reconstruction and physically constrained mineral abundance prediction
using a shared spectral representation.

\item \textbf{A scalable full-scene inference pipeline} for lunar mineral
mapping was developed using chunk-wise processing, abundance
reconstruction, Gaussian smoothing, and dominant-mineral assignment.

\item \textbf{An experimental evaluation} was performed on real IIRS
observations using spectral reconstruction metrics, abundance
consistency measures, and full-scene mineral distribution analysis.

\end{enumerate}



\section{Related Work}
\label{sec:related_work}

\subsection{Hyperspectral Analysis and Spectral Unmixing}

Hyperspectral analysis commonly employs spectral unmixing to decompose
mixed pixels into constituent endmembers and their fractional
abundances. Fully Constrained Least Squares Unmixing (FCLSU) enforces
non-negativity and abundance-sum constraints, whereas Non-negative
Matrix Factorization (NMF) estimates non-negative endmember and
abundance matrices through iterative optimization. However, FCLSU
depends on predefined endmembers and linear mixing assumptions, while
NMF is sensitive to initialization and may converge to local optima
\cite{bioucas2012unmixing,wang2023integration}. Nonlinear unmixing approaches address complex
material interactions but introduce additional modelling
complexity \cite{zhou2024blind}.

\subsection{Deep Learning and Transformer-Based Methods}

More recent
unmixing-oriented transformer architectures combine convolutional
feature extraction and attention to estimate endmembers and abundances
\cite{hadi2023dctnu}.

\begin{table}[t]
\centering
\caption{Comparison of representative hyperspectral analysis and
spectral unmixing approaches.}
\label{tab:related_work}
\scriptsize
\setlength{\tabcolsep}{2.5pt}
\renewcommand{\arraystretch}{1.15}
\resizebox{\columnwidth}{!}{%
\begin{tabular}{lcccccc}
\toprule
\textbf{Approach} &
\textbf{Endmember} &
\textbf{Abundance} &
\textbf{Recon.} &
\textbf{Supervision} &
\textbf{Complexity} &
\textbf{Task} \\
\midrule
FCLSU \cite{bioucas2012unmixing}
& Given & Yes & Yes & Unsupervised
& Iterative & Unmixing \\

NMF \cite{bioucas2012unmixing}
& Joint & Yes & Yes & Unsupervised
& Iterative & Unmixing \\

CNN-based HSI \cite{paoletti2019review}
& N/A & N/A & N/A & Supervised
& Architecture-dependent & Classification \\

DAEN \cite{su2019daen}
& Joint & Yes & Yes & Unsupervised
& Multi-stage & Unmixing \\

DCTNU \cite{hadi2023dctnu}
& Yes & Yes & Yes & Data-driven
& Multi-module & Unmixing \\

Double-DIP \cite{zhou2024blind}
& Joint & Yes & Yes & Unsupervised
& Multi-module & Unmixing \\

SpectralFormer \cite{hong2022spectralformer}
& N/A & N/A & N/A & Supervised
& Transformer & Classification \\

\textbf{Proposed method}
& \textbf{Predefined} & \textbf{Yes}
& \textbf{Yes} & \textbf{To be specified}
& \textbf{Transformer} & \textbf{Unmixing + mapping} \\
\bottomrule
\end{tabular}%
}
\end{table}

\subsection{Deep Learning-Based Spectral Unmixing}

Deep
Autoencoder Networks (DAEN) combine stacked autoencoders and variational
autoencoders for unsupervised endmember and abundance estimation
\cite{su2019daen}. Convolutional autoencoder-based methods further
incorporate spatial context into spectral reconstruction
\cite{yu2022multistage}.
Similarly, deep image prior-based approaches investigate blind
unmixing under linear and nonlinear mixing assumptions
\cite{zhou2024blind}.

\subsection{Research Gap}

The literature reveals three major limitations: (1) conventional
unmixing methods depend on predefined mixing assumptions or iterative
optimization; (2) most transformer-based HSI studies emphasize
classification, while deep unmixing methods commonly employ
autoencoder-based or spatial--spectral architectures; and (3)
computationally scalable one-dimensional spectral unmixing for
complete orbital observations remains relatively unexplored. Table~\ref{tab:related_work}
summarizes the methodological differences between representative
approaches.

%%%%%%%%%%%%%%%%%%%%%%%%%%%%%%%%%%%%%%%%%%%%%%%%%%%%%%%%%%%%%%%%%%%%%%

\section{Methodology}
\label{sec:methodology}

The proposed framework performs hyperspectral representation learning
and lunar mineral mapping using Chandrayaan-2 IIRS observations. As
shown in Fig.~\ref{fig:framework}, the complete pipeline comprises
spectral preprocessing, one-dimensional spectral tokenization, Vision
Transformer encoding, dual-task prediction, joint optimization, and
scalable full-scene inference.

%===========================================================
% Overall Architecture
%===========================================================
\begin{figure*}[!t]
    \centering
    \includegraphics[
        width=\textwidth,
        height=0.90\textheight,
        keepaspectratio
    ]{main.png}

    \caption{Overall architecture of the proposed one-dimensional
    Vision Transformer for hyperspectral spectral unmixing and lunar
    mineral mapping.}

    \label{fig:framework}
\end{figure*}
%===========================================================

\subsection{Data Preparation and Spectral Tokenisation}
\label{subsec:data_tokenisation}

The Chandrayaan-2 IIRS observations contain 256 spectral channels
spanning approximately $0.8$--$5.0~\si{\micro\metre}$. Each spatial
pixel is represented as a spectral signature
$\mathbf{x}\in\mathbb{R}^{C_0}$, where $C_0=256$. Noisy edge channels
and thermally contaminated wavelengths beyond
$3.5~\si{\micro\metre}$ were masked, retaining $C$ valid spectral
channels. The observations were obtained
from [VERIFY: exact IIRS data source and observation identifiers],
covering [VERIFY: geographical coordinates or study region].

Photometric correction was applied during preprocessing to improve
the radiometric consistency of IIRS observations
\cite{sharma2024photometric}. The remaining spectra were independently
normalized using min--max normalization:

\begin{equation}
x'_j =
\frac{x_j-\min(\mathbf{x})}
{\max(\mathbf{x})-\min(\mathbf{x})},
\qquad j=1,\ldots,C .
\label{eq:normalisation}
\end{equation}

The normalized observations were divided into training, validation,
and testing subsets using a [VERIFY: spatially disjoint or other]
partitioning strategy, with a split ratio of
[VERIFY: training/validation/testing percentages]. The corresponding
sample counts were [VERIFY: training pixels],
[VERIFY: validation pixels], and [VERIFY: testing pixels].
Spatial separation was maintained [VERIFY: explain the actual
procedure used to prevent spatial leakage].

Each normalized spectrum is divided into non-overlapping spectral
patches of length $P=16$:

\begin{equation}
\tilde{\mathbf{x}}^{(i)}\in\mathbb{R}^{P},
\qquad
N=\left\lfloor\frac{C}{P}\right\rfloor ,
\label{eq:patching}
\end{equation}

where $N$ denotes the number of spectral tokens. For each patch,
a linear projection maps the spectral values into a $D$-dimensional
embedding space, followed by a learnable positional embedding:

\begin{equation}
\mathbf{z}^{(i)}
=
\mathbf{W}_{E}\tilde{\mathbf{x}}^{(i)}
+
\mathbf{e}^{(i)} .
\label{eq:embedding}
\end{equation}

A learnable class token is prepended to aggregate global spectral
information:

\begin{equation}
\mathbf{Z}^{(0)}
=
[\mathbf{z}_{\mathrm{cls}},
\mathbf{z}^{(1)},\ldots,\mathbf{z}^{(N)}]
\in\mathbb{R}^{(N+1)\times D}.
\label{eq:input_sequence}
\end{equation}

The mineral categories
considered in this study comprise plagioclase, pyroxene, olivine,
ilmenite, spinel, glass, anorthosite, and mare basalt. Their
endmember/reference information was obtained from
[VERIFY: actual spectral library or reference source], with the
selection procedure described accordingly.

\subsection{One-Dimensional Vision Transformer Encoder}
\label{subsec:encoder}

The spectral token sequence is processed by a pre-layer
normalization ViT consisting of $L$ transformer blocks. Each block
contains multi-head self-attention (MHSA) and a feed-forward network
(FFN):

\begin{align}
\mathbf{Z}'_{\ell}
&=
\mathbf{Z}_{\ell-1}
+
\alpha\,\mathrm{MHSA}
\left(\mathrm{LN}(\mathbf{Z}_{\ell-1})\right),
\\
\mathbf{Z}_{\ell}
&=
\mathbf{Z}'_{\ell}
+
\alpha\,\mathrm{FFN}
\left(\mathrm{LN}(\mathbf{Z}'_{\ell})\right).
\label{eq:transformer_block}
\end{align}

Here, $\alpha$ denotes a learnable residual scaling parameter
that adjusts the contribution of each attention and feed-forward
branch to the residual representation.
Pre-Layer Normalization improves optimization stability
\cite{xiong2020layernorm}.

For attention head $h$,

\begin{equation}
\mathbf{Q}_{h}=\mathbf{Z}\mathbf{W}_{h}^{Q},
\qquad
\mathbf{K}_{h}=\mathbf{Z}\mathbf{W}_{h}^{K},
\qquad
\mathbf{V}_{h}=\mathbf{Z}\mathbf{W}_{h}^{V},
\end{equation}

and scaled dot-product attention is computed as

\begin{equation}
\mathrm{Attention}_{h}
=
\mathrm{softmax}
\left(
\frac{\mathbf{Q}_{h}\mathbf{K}_{h}^{\top}}
{\sqrt{d_k}}
\right)
\mathbf{V}_{h}.
\label{eq:attention}
\end{equation}

The FFN uses GELU activation \cite{hendrycks2016gelu}:

\begin{equation}
\mathrm{FFN}(\mathbf{z})
=
\mathbf{W}_{2}
\mathrm{GELU}
\left(
\mathbf{W}_{1}\mathbf{z}+\mathbf{b}_{1}
\right)
+\mathbf{b}_{2}.
\label{eq:ffn}
\end{equation}

The implemented encoder uses six transformer layers, a
256-dimensional embedding space, and eight attention heads.
The training configuration comprises a batch size of
[VERIFY], learning rate of [VERIFY], dropout rate of [VERIFY],
weight decay of [VERIFY], and [VERIFY] training epochs.
AdamW optimization with cosine learning-rate scheduling and
linear warm-up is employed. Experiments were conducted using
[VERIFY: processor/GPU, RAM, operating system, Python and
PyTorch versions].

\subsection{Dual-Task Learning and Optimization}
\label{subsec:dual_task}

The final class-token representation
$\mathbf{z}_{\mathrm{cls}}^{(L)}\in\mathbb{R}^{D}$ is shared by
two prediction heads. The reconstruction head estimates the
input spectrum:

\begin{equation}
\hat{\mathbf{x}}
=
\mathbf{W}_{R}\mathbf{z}_{\mathrm{cls}}^{(L)}
+\mathbf{b}_{R}.
\label{eq:reconstruction}
\end{equation}

The abundance head predicts the mineral contribution scores using
Softmax:

\begin{equation}
\hat{\mathbf{a}}
=
\mathrm{softmax}
\left(
\mathbf{W}_{A}\mathbf{z}_{\mathrm{cls}}^{(L)}
+\mathbf{b}_{A}
\right).
\label{eq:abundance}
\end{equation}

For $K=8$ mineral categories, the predicted abundance vector
satisfies the non-negativity and sum-to-one constraints:

\begin{equation}
\sum_{k=1}^{K}\hat{a}_{k}=1,
\qquad
\hat{a}_{k}\geq0.
\label{eq:abundance_constraint}
\end{equation}

The joint optimization objective is expressed as

\begin{equation}
\mathcal{L}
=
\mathcal{L}_{\mathrm{recon}}
+
\lambda\mathcal{L}_{\mathrm{abund}},
\label{eq:joint_loss}
\end{equation}

where $\mathcal{L}_{\mathrm{recon}}$ denotes the spectral
reconstruction loss, $\mathcal{L}_{\mathrm{abund}}$ represents
the abundance-learning objective, and $\lambda$ controls their
relative contribution. The reconstruction loss is calculated
using mean squared error:

\begin{equation}
\mathcal{L}_{\mathrm{recon}}
=
\frac{1}{BC}
\sum_{i=1}^{B}
\sum_{j=1}^{C}
(x_{ij}-\hat{x}_{ij})^2,
\label{eq:reconstruction_loss}
\end{equation}

where $B$ represents the batch size and $C$ denotes the number
of retained spectral channels.

[VERIFY: Specify the actual abundance-learning objective,
its reference targets, and the procedure used to generate
or obtain those targets. If independent abundance labels
are unavailable, clarify whether the objective uses
pseudo-labels or an unsupervised regularization strategy.]

The two objectives are jointly optimized. Dropout
is employed for regularization, and the model with the minimum
validation loss is retained for evaluation.

\subsection{Full-Scene Mineral Mapping}
\label{subsec:scene_mapping}

For complete orbital observations, the trained model is applied
using chunk-wise inference to control memory consumption. Each
scene is partitioned into spatial chunks of
[VERIFY: chunk height $\times$ chunk width] and processed
in mini-batches of [VERIFY: inference batch size].
Boundary regions are handled using [VERIFY: actual boundary
handling procedure] to preserve the original spatial
arrangement.

The predicted abundance vectors are reconstructed in their
original spatial arrangement. Gaussian smoothing is subsequently
applied independently to each abundance channel using a kernel
size of [VERIFY] and standard deviation of [VERIFY].
The dominant mineral at spatial position $(u,v)$ is determined
from the predicted abundance channels as

\begin{equation}
\mathcal{M}(u,v)
=
\arg\max_{k}\hat{a}_{k}(u,v).
\label{eq:mineral_map}
\end{equation}

The eight output channels are associated with the predefined
mineral categories in the order specified during model
configuration. Spatial continuity is examined through
[VERIFY: actual spatial validation procedure or qualitative
visual inspection].





\begin{algorithm}[t]
\caption{Full-Scene Lunar Mineral Mapping}
\label{alg:scene_mapping}
\begin{algorithmic}[1]
\Require IIRS hyperspectral cube $\mathbf{X}$, trained model
$f_{\theta}$, chunk size $C_L$, batch size $B_s$
\Ensure Mineral map $\mathcal{M}$ and abundance maps $\hat{\mathbf{A}}$

\State Partition $\mathbf{X}$ into spatial chunks of $C_L$ scan lines
\State Initialize abundance map $\hat{\mathbf{A}}$

\For{each spatial chunk $\mathbf{X}_c$}
    \State Extract valid pixel spectra from $\mathbf{X}_c$
    \State Normalize spectra using (\ref{eq:normalisation})

    \For{each mini-batch $\mathbf{X}_b$}
        \State $\hat{\mathbf{A}}_b \gets f_{\theta}(\mathbf{X}_b)$
    \EndFor

    \State Reconstruct spatial abundance arrangement for $\mathbf{X}_c$
    \State Insert $\hat{\mathbf{A}}_c$ into $\hat{\mathbf{A}}$
\EndFor

\State Apply Gaussian smoothing to each abundance channel

\For{each spatial position $(u,v)$}
    \State $\mathcal{M}(u,v)
    \gets \arg\max_k\hat{a}_k(u,v)$
\EndFor

\State \Return $\hat{\mathbf{A}},\mathcal{M}$
\end{algorithmic}
\end{algorithm} 

%%%%%%%%%%%%%%%%%%%%%%%%%%%%%%%%%%%%%%%%%%%%%%%%%%%%%%%%%%%%%%%%%%%%%%
\section{Results and Discussion}
\label{sec:results}


\subsection{Spectral Reconstruction}
\label{subsec:res_reconstruction}

Figure~\ref{fig:rmse_reconstruction} shows the representative best,
median, and worst spectral reconstruction cases. The reconstructed
spectra closely follow the corresponding IIRS observations and preserve
the major absorption characteristics near approximately $1.0$ and
$2.0~\si{\micro\metre}$.

\begin{figure}[t]
    \centering
    \includegraphics[width=0.85\columnwidth]{RMSEe1.png}
    \caption{Representative spectral reconstruction results for the
    best, median, and worst cases obtained by the proposed ViT.}
    \label{fig:rmse_reconstruction}
\end{figure}

The quantitative reconstruction results are presented in
Table~\ref{tab:quant_results}. The model achieved an RMSE of
$0.03637$, MAE of $0.01420$, and SAD of $0.02314$, indicating that the
reconstructed spectra retain the dominant characteristics of the input
observations. The corresponding SRE and ERGAS values are $28.78434$
and $31.14898$, respectively. These metrics assess spectral
reconstruction quality rather than the accuracy of the predicted
mineral abundances. Furthermore, the present results do not establish
superiority over conventional or learning-based unmixing methods;
comparative evaluation using the same data and metrics is required
to support such a conclusion.

\begin{table}[t]
    \centering
    \caption{Quantitative evaluation of the proposed ViT on the
    held-out IIRS test set.}
    \label{tab:quant_results}
    \footnotesize
    \setlength{\tabcolsep}{5pt}
    \renewcommand{\arraystretch}{1.10}
    \begin{tabular}{lc}
        \toprule
        \textbf{Metric} & \textbf{Proposed ViT} \\
        \midrule
        RMSE & 0.03637 \\
        MAE & 0.01420 \\
        SAD & 0.02314 \\
        SRE & 28.78434 \\
        ERGAS & 31.14898 \\
        PSNR (dB) & 28.78434 \\
        Abundance Sum Mean & 1.00000 \\
        Abundance Sum Std. & $4.35\times10^{-8}$ \\
        Abundance L1 Mean & 0.12500 \\
        Negative Abundance Ratio & 0.00000 \\
        Abundance Entropy Mean & 1.95543 \\
        \bottomrule
    \end{tabular}
\end{table}

\begin{table}[t]
    \centering
    \caption{Mineral-class distribution in the evaluated full-scene
    mineralogical map.}
    \label{tab:class_distribution}
    \footnotesize
    \setlength{\tabcolsep}{5pt}
    \renewcommand{\arraystretch}{1.10}
    \begin{tabular}{lrr}
        \toprule
        \textbf{Mineral Class} &
        \textbf{Pixel Count} &
        \textbf{Percentage (\%)} \\
        \midrule
        Plagioclase  & 43,948   & 3.19  \\
        Pyroxene     & 232,881  & 16.90 \\
        Olivine      & 6,797    & 0.49  \\
        Ilmenite     & 8,642    & 0.63  \\
        Spinel       & 205,677  & 14.92 \\
        Glass        & 0        & 0.00  \\
        Anorthosite  & 3,566    & 0.26  \\
        Mare Basalt  & 876,739  & 63.61 \\
        \midrule
        \textbf{Total} &
        \textbf{1,378,250} &
        \textbf{100.00} \\
        \bottomrule
    \end{tabular}
\end{table}

The learned latent representations shown in
Fig.~\ref{fig:learned_signatures} exhibit different patterns across
the principal mineral categories. However, latent
features should not be interpreted as physically measured mineral
endmember signatures or independent evidence of mineralogical
separability.

\begin{figure}[t]
    \centering
    \includegraphics[width=\columnwidth]{LESS.png}
    \caption{Learned latent spectral feature representations associated
    with the principal lunar mineral categories extracted by the
    proposed ViT encoder.}
    \label{fig:learned_signatures}
\end{figure}

\begin{figure}[t]
    \centering

    \begin{subfigure}[t]{0.47\columnwidth}
        \centering
        \includegraphics[width=\linewidth]{ADPE.png}
        \caption{Predicted abundance distributions.}
        \label{fig:abundance_dist}
    \end{subfigure}
    \hfill
    \begin{subfigure}[t]{0.47\columnwidth}
        \centering
        \includegraphics[width=\linewidth]{MMA.png}
        \caption{Mean predicted mineral abundance.}
        \label{fig:mean_abundance}
    \end{subfigure}

    \caption{Predicted mineral abundance distributions obtained by the
    proposed ViT.}
    \label{fig:abundance_analysis}
\end{figure}

\begin{figure}[t]
    \centering
    \includegraphics[width=\columnwidth,height=0.75\columnwidth,
    keepaspectratio]{osmm.png}
    \caption{Original Chandrayaan-2 IIRS observation and corresponding
    predicted mineral map generated by the proposed ViT.}
    \label{fig:spatial_map_comparison}
\end{figure}

\subsection{Mineral Abundance Estimation}
\label{subsec:res_abundance}

The predicted abundance distributions are shown in
Fig.~\ref{fig:abundance_dist}. Unlike hard classification, the proposed
framework produces continuous mineral fractions for each spectral
observation. This mathematical consistency is reflected by the abundance-sum
mean of $1.00000$ and standard deviation of $4.35\times10^{-8}$ reported
in Table~\ref{tab:quant_results}.

The mean abundance distribution in
Fig.~\ref{fig:mean_abundance} shows variation among the predicted
mineral classes, with mare basalt and spinel representing prominent
components of the evaluated observations.

The zero negative-abundance ratio reported in
Table~\ref{tab:quant_results} confirms that the predictions satisfy
the non-negativity constraint. Independent reference abundances or
mineralogical measurements are required to quantitatively validate
the predicted fractions.

\subsection{Full-Scene Mineral Mapping}
\label{subsec:res_spatial}

The trained model was applied to the complete Chandrayaan-2 IIRS
observations using the chunk-wise inference strategy described in
Section~\ref{sec:methodology}. The resulting predicted mineral map is
shown in Fig.~\ref{fig:spatial_map_comparison}.

The generated map exhibits spatially continuous mineral distributions
with reduced isolated pixel-level variations. However, spatial continuity alone does not confirm
geological correctness, and the mapped patterns require comparison
with independent lunar mineralogical references.

The scene-level distribution is presented in
Table~\ref{tab:class_distribution}. Mare basalt constituted the largest
mapped fraction ($63.61\%$), followed by pyroxene ($16.90\%$) and
spinel ($14.92\%$). Plagioclase accounted for $3.19\%$, whereas
olivine, ilmenite, and anorthosite represented smaller fractions.
No pixels were assigned to the glass class in the evaluated map.
These percentages represent the distribution of predicted dominant
mineral labels and do not indicate independently measured mineral
abundances.

%%%%%%%%%%%%%%%%%%%%%%%%%%%%%%%%%%%%%%%%%%%%%%%%%%%%%%%%%%%%%%%%%%%%%%
\section{Conclusion}
\label{sec:conclusion}

This paper presents a one-dimensional Vision Transformer framework for
hyperspectral unmixing and lunar mineral mapping using Chandrayaan-2 IIRS
observations. By representing spectra as spectral patches, the model uses
self-attention to capture long-range spectral dependencies and jointly
performs spectral reconstruction and physically constrained abundance
prediction through a dual-head architecture.

Current
limitations include predefined mineral endmembers, variability in degraded
orbital observations, and reliance on a single mission source. Future work
will explore adaptive endmember extraction, independent mineralogical
validation, cross-scene and cross-mission generalization, and physics-informed
self-supervised learning to improve the reliability and robustness of
planetary hyperspectral analysis.

\begin{thebibliography}{11}

\bibitem{chowdhury2020iirs}
Chowdhury, A.R., Chauhan, P., Bugiolacchi, R., et al.:
Planetary exploration of the Moon with Chandrayaan-2 Imaging Infrared Spectrometer (IIRS).
Current Science \textbf{118}(3), 368--375 (2020).
\doi{10.18520/cs/v118/i3/368-375}

\bibitem{chauhan2021hydration}
Chauhan, P., et al.:
Unambiguous detection of OH and H$_2$O on the Moon from Chandrayaan-2 Imaging Infrared Spectrometer reflectance data using the 3~$\mu$m hydration feature.
Current Science \textbf{121}(3), 391--401 (2021).
\doi{10.18520/cs/v121/i3/391-401}

\bibitem{dosovitskiy2021vit}
Dosovitskiy, A., Beyer, L., Kolesnikov, A., Weissenborn, D., Zhai, X.,
Unterthiner, T., Dehghani, M., Minderer, M., Heigold, G., Gelly, S.,
Uszkoreit, J., Houlsby, N.:
An Image is Worth 16$\times$16 Words: Transformers for Image Recognition at Scale.
In: International Conference on Learning Representations (ICLR) (2021).
\doi{10.48550/arXiv.2010.11929}

\bibitem{xiong2020layernorm}
Xiong, R., Yang, Y., He, D., Zheng, K., Zheng, S., Xing, C.,
Zhang, H., Lan, Y., Wang, L., Liu, T.:
On Layer Normalization in the Transformer Architecture.
In: International Conference on Machine Learning (ICML),
Proceedings of Machine Learning Research, vol.~119,
pp. 10524--10533 (2020).
\doi{10.48550/arXiv.2002.04745}

\bibitem{bioucas2012unmixing}
Bioucas-Dias, J.M., Plaza, A., Dobigeon, N., Parente, M., Du, Q.,
Gader, P., Chanussot, J.:
Hyperspectral Unmixing Overview: Geometrical, Statistical, and Sparse Regression-Based Approaches.
IEEE Journal of Selected Topics in Applied Earth Observations and Remote Sensing
\textbf{5}(2), 354--379 (2012).
\doi{10.1109/JSTARS.2012.2194696}

\bibitem{hendrycks2016gelu}
Hendrycks, D., Gimpel, K.:
Gaussian Error Linear Units (GELUs).
arXiv preprint arXiv:1606.08415 (2016).
\doi{10.48550/arXiv.1606.08415}

\bibitem{paoletti2019review}
Paoletti, M.E., Haut, J.M., Plaza, J., Plaza, A.:
Deep Learning Classifiers for Hyperspectral Imaging: A Review.
ISPRS Journal of Photogrammetry and Remote Sensing
\textbf{158}, 279--317 (2019).
\doi{10.1016/j.isprsjprs.2019.09.006}

\bibitem{hong2022spectralformer}
Hong, D., Han, Z., Yao, J., Gao, L., Zhang, B., Plaza, A.,
Chanussot, J.:
SpectralFormer: Rethinking Hyperspectral Image Classification with Transformers.
IEEE Transactions on Geoscience and Remote Sensing
\textbf{60}, 1--15 (2022).
\doi{10.1109/TGRS.2021.3130716}

\bibitem{he2021sst}
He, X., Chen, Y., Lin, Z.:
Spatial--Spectral Transformer for Hyperspectral Image Classification.
Remote Sensing \textbf{13}(3), 498 (2021).
\doi{10.3390/rs13030498}

\bibitem{sharma2024photometric}
Sharma, V., Chauhan, P., Kumar, A.:
Photometric Correction of Images Obtained from Chandrayaan-2 Imaging Infra-Red Spectrometer (IIRS) Data.
Advances in Space Research
\textbf{73}(4), 2105--2118 (2024).
\doi{10.1016/j.asr.2023.09.015}

\bibitem{zhang2016deeplearning}
Zhang, L., Zhang, L., Du, B.:
Deep Learning for Remote Sensing Data: A Technical Tutorial on the State of the Art.
IEEE Geoscience and Remote Sensing Magazine
\textbf{4}(2), 22--40 (2016).
\doi{10.1109/MGRS.2016.2540798}

\bibitem{shu2024datn}
Shu, Z., Wang, Y., Yu, Z.:
Dual Attention Transformer Network for Hyperspectral Image Classification.
Engineering Applications of Artificial Intelligence
\textbf{127}, 107351 (2024).


\bibitem{su2019daen}
Su, Y., Li, J., Plaza, A., Marinoni, A., Gamba, P., Chakravortty, S.:
DAEN: Deep Autoencoder Networks for Hyperspectral Unmixing.
IEEE Transactions on Geoscience and Remote Sensing
\textbf{57}(7), 4309--4321 (2019).
\doi{10.1109/TGRS.2018.2890633}

\bibitem{wang2023integration}
Chen, J., Zhao, M., Wang, X., Richard, C., Rahardja, S.:
Integration of Physics-Based and Data-Driven Models for Hyperspectral
Image Unmixing: A Summary of Current Methods.
IEEE Signal Processing Magazine
\textbf{40}(2), 61--74 (2023).
\doi{10.1109/MSP.2022.3208987}

\bibitem{hadi2023dctnu}
Hadi, F., Yang, J., Farooque, G., Xiao, L.:
Deep Convolutional Transformer Network for Hyperspectral Unmixing.
European Journal of Remote Sensing
\textbf{56}(1), Article 2268820 (2023).
\doi{10.1080/22797254.2023.2268820}

\bibitem{zhou2024blind}
Zhou, C., Rodrigues, M.R.D.:
Hyperspectral Blind Unmixing Using a Double Deep Image Prior.
IEEE Transactions on Neural Networks and Learning Systems
\textbf{35}(11), 16478--16492 (2024).
\doi{10.1109/TNNLS.2023.3294714}

\bibitem{yu2022multistage}
Yu, Y., Ma, Y., Mei, X., Fan, F., Huang, J., Li, H.:
Multi-Stage Convolutional Autoencoder Network for Hyperspectral Unmixing.
International Journal of Applied Earth Observation and Geoinformation
\textbf{113}, Article 102981 (2022).
\doi{10.1016/j.jag.2022.102981}

\end{thebibliography}
\end{document}
